\subsection{Graph Reduction: Definitions}
\label{sec:prelim:reduction}
A reduction operator $R$ maps a graph $G$ to a smaller graph $R(G)$ standing in for it. The
taxonomy follows \citet{hashemi2024survey}, with one substitution: partitioning replaces
condensation as the third family, since partitioning is what a distributed setup forces on
a graph too large for one device (\ref{sec:relwork:scaling:distributed}), and condensation
is excluded outright (\ref{sec:relwork:reduction:summarization}).

% A reduction operator $R$ maps a graph $G$ to a smaller graph $R(G)$ intended to stand in
% for it. The taxonomy used throughout follows the survey of \citet{hashemi2024survey},
% which divides the field into sparsification, coarsening and condensation, with one
% substitution: partitioning replaces condensation as the third family. Partitioning is
% what a distributed training setup forces on a graph that does not fit one device
% (\ref{sec:relwork:scaling:distributed}), whereas condensation is excluded from this study
% outright, for the reasons given in~\ref{sec:relwork:reduction:summarization}. The
% substitution is deliberate, not an error of transcription.

\subsubsection{Partitioning}
\label{sec:prelim:reduction:partitioning}
Partitioning splits $\mathcal{V}$ into $k$ disjoint blocks and deletes every edge whose
endpoints fall in different blocks, so every node survives and whatever it compresses is the
edge set alone.

% Partitioning splits $\mathcal{V}$ into $k$ disjoint blocks and deletes every edge whose
% endpoints fall in different blocks. It is the one family that keeps every node: only
% spanning edges are lost, so the node count of $R(G)$ always equals that of $G$ and any
% compression it achieves is compression of the edge set alone. What varies between
% partitioners is the objective by which blocks are chosen
% (\ref{sec:method:reduction:partitioning}).

\subsubsection{Sparsification}
\label{sec:prelim:reduction:sparsification}
Sparsification removes edges, or nodes with their incident edges, while preserving some
structural property of the original, such as pairwise distances up to a stretch factor.
Edge-removing and node-removing variants differ in whether $|\mathcal{V}|$ moves at all,
forcing the two-ratio convention of~\ref{sec:prelim:reduction:measuring}.

% Sparsification removes edges, or nodes together with their incident edges, while
% preserving some structural property of the original, such as pairwise distances up to a
% stretch factor. The two variants are not interchangeable for measurement: an edge-removing
% method leaves $|\mathcal{V}|$ untouched while a node-removing method shrinks it, and the
% distinction is what forces the two-ratio convention
% of~\ref{sec:prelim:reduction:measuring}.

\subsubsection{Summarization and Coarsening}
\label{sec:prelim:reduction:summarization}
Summarization merges nodes into \emph{super-nodes} by a merge map, a surjection
$\mu \colon \mathcal{V} \rightarrow \mathcal{V}'$ onto the cluster set, and the resulting
quotient graph connects two clusters whenever any of their members were connected. The
number of member edges a super-edge stands for is its multiplicity,
\begin{equation}
    w_{u'v'} = \bigl| \{\, (u, v) \in \mathcal{E} :
        \mu(u) = u',\ \mu(v) = v' \,\} \bigr|,
    \label{eq:prelim:multiplicity}
\end{equation}
and a quotient that keeps attribute-distinct super-edges parallel is a multigraph. Because
$\mu$ is a function, a coarsening \emph{is} a partition of $\mathcal{V}$, so a tolerance
band such as ``within one level'' cannot define one: the relation it induces is not
transitive (\ref{sec:method:summary:domain}).

% Summarization merges nodes into \emph{super-nodes}. A coarsening is specified by a merge
% map, a surjection $\mu \colon \mathcal{V} \rightarrow \mathcal{V}'$ onto the cluster set,
% and the resulting quotient graph connects two clusters whenever any of their members were
% connected. The number of member edges a super-edge stands for is its multiplicity,
% and when super-edges differing in an attribute are kept parallel rather than collapsed,
% the quotient is a multigraph. Because $\mu$ is a function, its preimages partition
% $\mathcal{V}$: a coarsening \emph{is} a partition of the node set. A tolerance band such
% as ``within one level'' therefore cannot define one, since the relation it induces is not
% transitive, a fact the domain-specific candidate of~\ref{sec:method:summary:domain} has
% to work around.

\subsubsection{Exact Summarization}
\label{sec:prelim:reduction:exact}
A coarsening is \emph{exact} for an encoder when the encoder cannot tell the quotient from
the original. Writing $f_\theta$ for an $L$-layer instance of~\eqref{eq:prelim:mp} with a
graph-level readout, and $R$ for the coarsening with merge map $\mu$, exactness is
\begin{equation}
    f_\theta \bigl( R(G) \bigr) = f_\theta(G)
    \quad \text{for every } \theta .
    \label{eq:prelim:exact}
\end{equation}
Quantifying over every $\theta$, not just the weights that were trained, separates exactness
from an approximation that merely happens to be accurate. Its conditions
\citep{bollen2023exact} constrain three separate objects, and a method meeting two of the
three is not exact.

\paragraph{(a) The quotient.}
Merged nodes must be indistinguishable to $L$ rounds of colour refinement
(\ref{sec:prelim:gnn:expressivity}),
\begin{equation}
    \mu(u) = \mu(v) \;\Longrightarrow\; c_u^{(L)} = c_v^{(L)} ,
    \label{eq:prelim:exact-quotient}
\end{equation}
so $\mu$'s partition refines the $L$-round colour partition, the coarsest admissible merge.
Two side conditions apply: the initial colouring $c^{(0)}$ must separate distinct feature
rows, and~\eqref{eq:prelim:cr} must run over the same neighbourhood and edge attributes the
encoder aggregates over, since refining by fan-out while aggregating over fan-in merges
nodes the encoder can still tell apart. The count cap must likewise match what the
aggregation can count: capping below that makes the colour a coarser invariant than the
aggregate computed, which is why bisimulation is exact for a set-valued aggregation and
lossy for a counting one.

\paragraph{(b) The aggregation.}
Collapsing $w_{u'v'}$ parallel edges into one super-edge replaces $w_{u'v'}$ identical
messages by a single message, so $\bigoplus$ must read the multiplicity
of~\eqref{eq:prelim:multiplicity}, not the set of distinct neighbours. For the summing
encoder used here, that multiplicity enters as a coefficient on the message:
\begin{equation}
    \bigoplus_{u \in \mathcal{N}(v)} \psi \bigl( h_u, e_{uv} \bigr)
    \;=\; \sum_{u' \in \mathcal{N}'(\mu(v))}
        w_{u' \mu(v)} \, \psi \bigl( h_{u'}, e_{u' \mu(v)} \bigr) ,
    \label{eq:prelim:exact-aggregation}
\end{equation}
with $\mathcal{N}'$ the quotient neighbourhood. Where the multiplicity enters decides
whether~\eqref{eq:prelim:exact-aggregation} holds: $\psi$ is nonlinear, so
$\psi(h, w_{u'v'} e) \neq w_{u'v'} \, \psi(h, e)$, and a multiplicity carried as an edge
attribute breaks the equality a coefficient outside $\psi$ preserves
(\ref{sec:method:architecture:exact}).

\paragraph{(c) The readout.}
Conditions (a) and (b) give $h_{\mu(v)}^{l} = h_v^{l}$ at every layer, so pooling over
$\mathcal{V}'$ must weight each super-node by its member count $m_{v'} = | \mu^{-1}(v') |$
to reproduce pooling over $\mathcal{V}$:
\begin{equation}
    \frac{1}{\left| \mathcal{V} \right|} \sum_{v \in \mathcal{V}} h_v^{L}
    \;=\; \frac{\sum_{v' \in \mathcal{V}'} m_{v'} \, h_{v'}^{L}}
               {\sum_{v' \in \mathcal{V}'} m_{v'}} .
    \label{eq:prelim:exact-readout}
\end{equation}
An unweighted mean violates~\eqref{eq:prelim:exact-readout}, giving a super-node of fifty
members the weight of one. The member counts a summarized feature schema already carries
make the weighting available.

Colour refinement at unbounded count cap is the one reduction here that satisfies all three
(\ref{sec:method:summary:wl}). Sparsification and partitioning produce no quotient and
delete edges the encoder aggregates over, so~\eqref{eq:prelim:exact} fails for them by
construction: exactness is unavailable to those families, not merely unachieved.

\subsubsection{Measuring Reduction: Compression Ratio}
\label{sec:prelim:reduction:measuring}
Reduction is reported as retention, kept over original, separately per node and per edge:
\begin{equation}
    \eta_{\mathcal{V}}(G) =
        \frac{\left| \mathcal{V}(R(G)) \right|}{\left| \mathcal{V}(G) \right|},
    \qquad
    \eta_{\mathcal{E}}(G) =
        \frac{\left| \mathcal{E}(R(G)) \right|}{\left| \mathcal{E}(G) \right|},
    \label{eq:prelim:retention}
\end{equation}
with reduction $1 - \eta$. One number cannot carry both: partitioning and edge-removing
sparsification hold $\eta_{\mathcal{V}} = 1$ while lowering $\eta_{\mathcal{E}}$, and
node-removing methods move the two together, measured on this corpus and matched across
methods in~\ref{sec:method:experiment:metrics:reduction}.

% Reduction is reported as retention, kept over original, separately per node and per edge:
% with reduction $1 - \eta$. One number cannot carry both: partitioning and edge-removing
% sparsification hold $\eta_{\mathcal{V}} = 1$ while lowering $\eta_{\mathcal{E}}$, and
% node-removing methods move the two together. Both ratios are therefore reported
% everywhere, and never collapsed (\ref{sec:method:experiment:metrics:reduction}). Methods
% further differ in whether compression is a knob or an outcome: some accept a target
% ratio, while for others $\eta$ is a fixed property of the input graph. Comparing methods
% at matched compression is consequently a design problem rather than a reporting
% convention, and~\ref{sec:method:summary:random} is its solution.
